% Image credits for the photographs and AI illustrations of Book 4
% (University Chemistry, Year 3). Machine-readable license records live in
% images/book4/CREDITS.md; keep the two files in sync. The Book 4 writer
% owns this file: add one \item per photograph (in an itemize, after the
% first paragraph) as photographs land.
\chapter*{Créditos das imagens}

Os desenhos deste livro são originais. As fotografias, quando usadas, são
utilizadas sob as respectivas licenças livres, com agradecimento aos seus
autores; os links completos para as fontes e as URLs das licenças de cada
fotografia estão listados em \texttt{images/book4/CREDITS.md} no
repositório do livro.
Os pictogramas de perigo são os do Sistema Globalmente Harmonizado de
Classificação e Rotulagem de Produtos Químicos, tal como desenhados para o
pacote \texttt{ghsystem} do \TeX\ Live.

\bigskip
Ao lado das fotografias e dos desenhos originais, algumas figuras são
ilustrações geradas por IA, produzidas com a geração de imagens da OpenAI a
partir de instruções escritas pelos editores do livro, e revisadas quanto à
exatidão química antes de serem incluídas. A instrução completa de cada
imagem gerada está registrada em \texttt{images/book4/ai/PROMPTS.md} no
repositório.
\bigskip
\noindent Fotografias, por capítulo:
\begin{itemize}\small
  \item Cap.~1: Erwin Schr\"odinger, 1933, Fundação Nobel, domínio público
        (Wikimedia Commons).
  \item Cap.~4: cristais de neve, Wilson A. Bentley (por volta de 1902), domínio público
        (Wikimedia Commons).
  \item Cap.~6: antenas de radiotelescópio em um planalto elevado, ESO/B.~Tafreshi
        (twanight.org), CC BY 4.0 (Wikimedia Commons).
  \item Cap.~7: fluorita à luz do dia e sob luz ultravioleta, de Masha
        Milshina, CC BY 4.0 (Wikimedia Commons).
  \item Cap.~8: um laboratório de RMN, de Shandchem, CC BY 2.0 (Wikimedia
        Commons).
  \item Cap.~9: William Lawrence Bragg, Fundação Nobel, domínio público
        (Wikimedia Commons).
  \item Cap.~10: o túmulo de Ludwig Boltzmann, de Feldkurat Katz,
        CC BY-SA 4.0 (Wikimedia Commons).
  \item Cap.~13: espirais de Belousov--Zhabotinsky, de Stephen Morris, CC BY 2.0
        (Wikimedia Commons).
  \item Cap.~14: o buraco na camada de ozônio sobre a Antártida, NASA, domínio público (Wikimedia
        Commons).
  \item Cap.~16: micrografia eletrônica da colmeia de um conversor catalítico, de
        Galadrid, CC BY-SA 4.0 (Wikimedia Commons).
  \item Cap.~17: raios de sol e o efeito Tyndall, de Kevin c higgins, CC BY-SA 4.0
        (Wikimedia Commons).
  \item Cap.~18: cristais de rubi e de esmeralda, de Robert M. Lavinsky, CC BY-SA 3.0
        (Wikimedia Commons).
  \item Cap.~20: cristais de ferroceno, de Andrea Sella, CC BY 4.0 (Wikimedia
        Commons).
  \item Cap.~22: um lingote de silício monocristalino, de Sebastian Wallroth, domínio
        público (Wikimedia Commons).
  \item Cap.~23: uma flor sobre aerogel de sílica, NASA/JPL, domínio público; a
        taça de Licurgo, de Marie-Lan Nguyen, CC BY 2.5 (Wikimedia Commons).
  \item Cap.~24: um límulo do Atlântico, de Kaldari, CC0 (Wikimedia
        Commons).
  \item Cap.~27: flores de piretro, de Soramimi, CC BY-SA 4.0 (Wikimedia
        Commons).
  \item Cap.~30: casca de um teixo do Pacífico, de JOE BLOWE (theforestprimeval),
        CC BY-SA 2.0 (Wikimedia Commons).
  \item Cap.~32: uma floração de algas vista pelo satélite Landsat 8, USGS e NASA,
        domínio público.
  \item Cap.~33: uma linha dupla de vácuo e de gás inerte, de Polimerek, CC BY-SA
        3.0 (Wikimedia Commons).
\end{itemize}
\clearpage
